Table~\ref{tab:related} compares the main analytical focus and supporting evidence of selected related studies with the present framework. The comparison describes complementary study designs and does not establish comparative performance.

\begingroup
\small
\setlength{\tabcolsep}{4pt}
\renewcommand{\arraystretch}{1.15}
\begin{longtable}{@{}>{\raggedright\arraybackslash}p{0.16\linewidth}>{\raggedright\arraybackslash}p{0.27\linewidth}>{\raggedright\arraybackslash}p{0.27\linewidth}>{\raggedright\arraybackslash}p{0.23\linewidth}@{}}
\caption{Selected liver transcriptomic studies and the present study}\label{tab:related}\\
\toprule
Study and data & Reported analytical focus & Biological context and supporting evidence & Relationship to the present framework \\
\midrule
\endfirsthead
\toprule
Study and data & Reported analytical focus & Biological context and supporting evidence & Relationship to the present framework \\
\midrule
\endhead
\bottomrule
\endfoot
Hoang et al., 2019 \cite{ref4}; 78 biopsies & Gene-expression associations across the histological spectrum and gene-level severity profiles. & Disease-associated molecular heterogeneity; longitudinal follow-up identified as a validation need. & Establishes a precedent for stage-related expression analysis; the present study combines evidence across cohorts. \\
\addlinespace
Govaere et al., 2020 \cite{govaere2020}; 206 NAFLD discovery cases & Transcription across the NAFLD spectrum and a 25-gene signature associated with fibrosing steatohepatitis. & Independent replication cohort, serum-protein follow-up and selected functional experiments. & Demonstrates stage-associated signatures and validation; the present study emphasises five-stage gene profiles and explicit prioritisation scores. \\
\addlinespace
Piras and DiStefano, 2024 \cite{piras2024}; ten datasets, 1,058 samples & Random-effects meta-analysis comparing normal liver, MASLD and MASH. & Genetic-association integration, coexpression and regulatory networks, and module-preservation assessment. & Provides cross-cohort disease-group evidence; the present framework separately analyses ordered fibrosis stages and disease groups. \\
\addlinespace
Yang et al., 2026 \cite{yang2026}; cohorts across liver disease aetiologies & Shared and aetiology-associated fibrogenic expression programmes and hub-gene prioritisation. & Network and enrichment analyses; RT-qPCR in ten fibrotic and ten non-fibrotic tissues and selected protein measurements. & Addresses conservation across aetiologies; the present framework focuses on stage-associated patterns and curated graph annotations. \\
\addlinespace
Maiers et al., 2026 \cite{maiers2026}; approximately 1,000 patients & Advanced versus mild or absent fibrosis (F3--F4 versus F0--F1), excluding F2 from this comparison. & Pathway, cell-type and ligand--receptor analyses; partial confirmation of selected genes in an independent cohort. & Characterises advanced fibrosis; the present framework retains F0--F4 profiles and evaluates four adjacent-stage contrasts. \\
\midrule
\textbf{Present study}; eight datasets, 1,085 samples; 610 staged patient-level observations & Separate ordered fibrosis and disease-group analyses; five gene-profile clusters and four adjacent-stage contrasts. & OptimusKG integration, multi-evidence prioritisation of 447 Tier 1 genes and 23 shortlisted candidates, and leave-one-study-out assessment. & Integrates these evidence types within one framework. Robustness assessment is internal; longitudinal, external and functional validation remain necessary. \\
\end{longtable}
\endgroup

